% --- Perfil ---
\section{Perfil}
\begin{onecolentry}
	Soy graduado en Ingeniería de Telecomunicación, habiendo cursado el grado y el máster en la Universidad de Alcalá. Dentro del campo de la ingeniería, estoy especialmente interesado en los sistemas de radiofrecuencia y en el diseño de sistemas de comunicaciones, habiendo orientado mi trayectoria profesional hacia la electrónica de RF, la instrumentación y la integración de subsistemas para adquisición, medida y validación. He adquirido experiencia en entornos en los que convergen hardware, comunicaciones e instrumentación avanzada, trabajando con plataformas de altas prestaciones y con procesos de caracterización y evaluación técnica, lo que me ha permitido desarrollar una visión aplicada del diseño, la implementación y la validación de soluciones dentro del ámbito de los sistemas de RF.
\end{onecolentry}

\vspace{-5pt}

% --- Educación ---
\section{Educación}

\begin{onecolentry}
	\begin{tabularx}{\textwidth}{@{}>{\raggedright\arraybackslash}X>{\raggedright\arraybackslash}p{3.6cm}>{\raggedleft\arraybackslash}p{2.4cm}@{}}
		\textbf{Máster en Ingeniería en Tecnologías de Telecomunicación} & Universidad de Alcalá & \textit{Febrero 2026}
	\end{tabularx}
	
	\vspace{0.05 cm}
	\hspace{0.2cm}Premio extraordinario
\end{onecolentry}

\vspace{0.2 cm}

\begin{onecolentry}
	\begin{tabularx}{\textwidth}{@{}>{\raggedright\arraybackslash}X>{\raggedright\arraybackslash}p{3.6cm}>{\raggedleft\arraybackslash}p{2.4cm}@{}}
		\textbf{Grado en Ingeniería en Tecnologías de Telecomunicación} & Universidad de Alcalá & \textit{Julio 2024}
	\end{tabularx}
	
	\vspace{0.05 cm}
	\hspace{0.2cm}2.º puesto de promoción y premio extraordinario
\end{onecolentry}

\vspace{-7pt}

% --- Experiencia ---
\section{Experiencia}

\begin{twocolentry}{\textit{Septiembre 2024 – Enero 2026}}
	\textbf{Plataformas técnicas y validación} \\
	Universidad de Alcalá, Alcalá de Henares
\end{twocolentry}
\begin{onecolentry}
	\begin{highlights}
		\item Diseñé y desplegué entornos reproducibles sobre contenedores y orquestadores para pruebas, experimentación y validación técnica.
		\item Desarrollé herramientas en Python, MATLAB y C para adquisición, procesado y análisis de datos en plataformas de altas prestaciones.
		\item Integré flujos automatizados de despliegue y validación para facilitar la repetibilidad experimental y la integración de componentes software.
	\end{highlights}
\end{onecolentry}

\vspace{0.2 cm}

\begin{twocolentry}{\textit{Febrero 2024 – Agosto 2024}}
	\textbf{Sistemas, redes y soporte corporativo} \\
	Correos y Telégrafos S.A., Madrid
\end{twocolentry}
\begin{onecolentry}
	\begin{highlights}
		\item Trabajé en la operación y documentación de infraestructuras y servicios técnicos en entorno empresarial.
		\item Gestioné servicios de seguridad y autenticación y participé en la administración de sistemas Linux empresariales.
		\item Colaboré en despliegues y automatizaciones ligados a la operación de plataformas corporativas.
	\end{highlights}
\end{onecolentry}

\vspace{0.2 cm}

\begin{twocolentry}{\textit{Enero 2023 – Enero 2024}}
	\textbf{Ingeniería e investigación aplicada} \\
	Universidad de Alcalá / Correos y Telégrafos S.A., Alcalá de Henares
\end{twocolentry}
\begin{onecolentry}
	\begin{highlights}
		\item Desarrollé plataformas de análisis de datos y soporte experimental para proyectos de investigación.
		\item Implementé soluciones de computación multihilo y en clúster para tareas de procesado intensivo y validación técnica.
		\item Me encargué de la administración de entornos Linux y colaboré en actividades de investigación aplicada y soporte experimental.
		\item Participé como coautor en publicaciones científicas \href{https://orcid.org/0009-0000-6862-1872}{ORCID}.
	\end{highlights}
\end{onecolentry}

\vspace{-7pt}

% --- Aptitudes ---
\section{Aptitudes}

\begin{onecolentry}
	\textbf{Programación:} C/C++, Python, MATLAB, VHDL, Bash
\end{onecolentry}
\begin{onecolentry}
	\textbf{Instrumentación y CAD:} LabVIEW, PSpice, KiCad, AutoCAD, Autodesk Inventor
\end{onecolentry}
\begin{onecolentry}
	\textbf{MCUs:} Xilinx Zynq, STM32, ESP32, NXP 17xx, Renesas Synergy
\end{onecolentry}
\begin{onecolentry}
	\textbf{Sistemas de RF:} Electrónica de RF, diseño de sistemas de comunicaciones, instrumentación de radiofrecuencia, caracterización y validación de subsistemas RF
\end{onecolentry}